% ============================================================================
\section{模曲线 $Y(\Gamma)$ 的复结构}
\label{sec:modular-curve-complex-structure}
\label{sec:y-gamma-complex-structure}
% ============================================================================

\textbf{我们到底想解决什么问题？}

在\cref{sec:modular-curves-moduli,sec:congruence-subgroups} 中，我们已经把
$Y(\Gamma)=\Gamma\backslash\HH$ 视为一个分类空间：它的点代表带有某种级结构的椭圆曲线。
但如果它只是一个"点的集合"，那我们几乎什么都做不了：
不能自然地谈全纯函数，不能谈亚纯函数，也不能谈微分形式与线丛。

\textbf{核心想法非常朴素。}
上半平面 $\HH$ 本身已经是黎曼面，
所以我们希望把商映射
\[
  \pi\colon \HH \longrightarrow Y(\Gamma)=\Gamma\backslash\HH
\]
附近的局部形状读出来：
在绝大多数点，取商不会造成任何扭曲；
只有在稳定子非平凡的点——也就是椭圆点——附近，商空间会像
"把几个扇形粘在一起"，这时就必须换一个更合适的局部坐标。

\textbf{一旦这件事做成，会得到什么？}
答案是：$Y(\Gamma)$ 成为一个非紧黎曼面。
于是模形式可以被理解为这个黎曼面上的几何对象；
再往前走，紧化后的 $X(\Gamma)$ 有亏格，
而亏格将直接出现在\cref{ch:dimension} 的维数公式里。

\begin{proposition}[普通点与椭圆点的局部模型]
\label{prop:y-gamma-local-model}
设 $\Gamma\subset \SL_2(\Z)$ 为有限指数子群，$\tau_0\in\HH$。
记 $\Gamma_{\tau_0}=\operatorname{Stab}_\Gamma(\tau_0)$。
则有两种情形：
\begin{enumerate}
  \item 若 $\Gamma_{\tau_0}=\{\pm I\}$（在 $\HH$ 上即平凡稳定子），
    则存在 $\tau_0$ 的小邻域 $U$，使得 $\pi|_U\colon U\to \pi(U)$ 是同胚。
    因而 $\tau-\tau_0$ 给出 $Y(\Gamma)$ 在 $[\tau_0]$ 附近的局部坐标。
  \item 若 $\Gamma_{\tau_0}$ 的像在 $\PSL_2(\R)$ 中阶为 $e>1$，
    则 $[\tau_0]$ 是一个阶 $e$ 的椭圆点。
    在某个以 $\tau_0$ 为中心的局部坐标 $u$ 下，稳定子作用可写成
    $u\mapsto \zeta_e u$，其中 $\zeta_e=e^{2\pi i/e}$。
    因而商空间上的自然局部坐标是
    \[
      w=u^e.
    \]
    直观上，这可以理解为 $w=(\tau-\tau_0)^e$。
\end{enumerate}
\end{proposition}

\begin{proof}[思路]
$\Gamma$ 在 $\HH$ 上作用不连续，因此若稳定子平凡，就可以取一个足够小的圆盘 $U$，
使得除恒等元外没有任何 $\gamma\in\Gamma$ 把 $U$ 与自身相交。
这时商映射在 $U$ 上没有折叠，局部坐标当然仍可直接用 $\tau$。

真正的新现象出现在椭圆点。
由有限阶 Möbius 变换的局部线性化，
可取坐标
\[
  u=\frac{\tau-\tau_0}{\tau-\overline{\tau_0}},
\]
把 $\tau_0$ 送到 $0$，并把附近小邻域送到单位圆盘中的小圆盘；
此时稳定子生成元在该坐标下就是旋转 $u\mapsto \zeta_e u$。
商掉这个旋转后，唯一不变而又能区分轨道的坐标正是 $w=u^e$。
所以在椭圆点处，$u$ 不是商空间坐标，但 $u^e$ 是。
\end{proof}

% figures/ch02-elliptic-local-coordinate.tex
% 椭圆点附近的局部坐标：u-map on the disk, quotient by rotation, w=u^e
\begin{figure}[ht]
\centering
\begin{tikzpicture}[>=Stealth, font=\small, scale=1.05]
  % left disk
  \draw[thick] (-4,0) circle (1.7);
  \fill[red!80!black] (-4,0) circle (2pt);
  \node[below] at (-4,-1.9) {上半平面中的小邻域 $U$};
  \node[left, red!80!black] at (-4.05,0.05) {$\tau_0$};

  % spokes / fan
  \foreach \ang in {15,87,159,231,303} {
    \draw[blue!70!black, thick, ->] (-4,0) -- ++({1.35*cos(\ang)},{1.35*sin(\ang)});
  }
  \node[blue!70!black, align=center] at (-4,2.35)
    {$e$ 个方向在商空间中\\会被识别在一起};
  \node[align=center] at (-4,-2.55)
    {$u$ 是以 $\tau_0$ 为中心的局部坐标\\直观上可把 $u$ 想成 $\tau-\tau_0$};

  % middle arrow
  \draw[->, very thick] (-2.1,0) -- (-0.6,0);
  \node[above] at (-1.35,0.15) {按稳定子取商};

  % middle disk with identified rays
  \draw[thick] (1.2,0) circle (1.7);
  \fill[red!80!black] (1.2,0) circle (2pt);
  \draw[orange!80!black, line width=1.1pt, ->] (1.2,0) -- ++(1.35,0);
  \node[below] at (1.2,-1.9) {$U/\operatorname{Stab}_\Gamma(\tau_0)$};
  \node[orange!80!black, align=center] at (1.2,2.3)
    {$u \sim \zeta_e u$\\（旋转 $2\pi/e$ 后是同一点）};

  % right arrow
  \draw[->, very thick] (3.25,0) -- (4.75,0);
  \node[above] at (4.0,0.15) {$w=u^e$};

  % right plane patch
  \draw[thick, rounded corners=10pt] (5.4,-1.35) rectangle (8.6,1.35);
  \fill[red!80!black] (7.0,0) circle (2pt);
  \draw[green!60!black, thick, ->] (7.0,0) -- ++(1.1,0.55);
  \node[below] at (7.0,-1.65) {真正的局部坐标盘};
  \node[green!60!black, align=center] at (7.0,2.2)
    {$w$ 把 $e$ 个方向压成一个方向\\所以商空间重新变得像 $\C$};
\end{tikzpicture}
\caption{椭圆点附近的局部模型。若 $\tau_0$ 的稳定子阶为 $e$，则在适当坐标 $u$ 下稳定子作用是 $u \mapsto \zeta_e u$，因此商空间的自然坐标是 $w=u^e$。直观上可以把它理解为 $w=(\tau-\tau_0)^e$。}
\label{fig:elliptic-local-coordinate}
\end{figure}


\textbf{为什么要在椭圆点处改坐标？}
如果还硬把 $\tau-\tau_0$ 当作商空间的坐标，
那么 $u$ 与 $\zeta_eu$ 会代表同一个商点，
这就破坏了"一个坐标值对应一个点"这件事。
换成 $w=u^e$ 之后，旋转的 $e$ 个方向都被压成同一个方向，
商空间又重新看起来像一块普通的复平面。

\begin{definition}[$Y(\Gamma)$ 的复图册]
\label{def:y-gamma-atlas}
$Y(\Gamma)=\Gamma\backslash\HH$ 上的\textbf{复图册}定义如下：
\begin{itemize}
  \item 在普通点 $[\tau_0]$ 附近，取上面命题中的小邻域 $U$，
    用 $\tau-\tau_0$ 作为局部坐标；
  \item 在阶 $e$ 的椭圆点 $[\tau_0]$ 附近，取线性化坐标 $u$，
    用 $w=u^e$ 作为局部坐标。
\end{itemize}
由这些坐标卡生成的极大图册，称为 $Y(\Gamma)$ 的\textbf{复结构}。
\end{definition}

\begin{proposition}
\label{prop:y-gamma-riemann-surface}
对任意有限指数子群 $\Gamma\subset\SL_2(\Z)$，商空间 $Y(\Gamma)$ 在上述图册下是一个非紧黎曼面。
\end{proposition}

\begin{proof}
覆盖性来自于每个点都属于普通点或椭圆点两类之一。
坐标变换之所以全纯，是因为它们都来自 $\HH$ 上的全纯坐标变换：
普通点之间如此，普通点与椭圆点之间也是如此，
只是椭圆点处先经过 $u\mapsto u^e$ 再取逆。
因此得到一个合法的复图册。

它是非紧的，因为有限指数子群都有尖点；
换句话说，总可以沿着某个尖点方向让 $\im\tau\to\infty$ 而逃出任何紧集。
这些"缺失的无穷远点"将在下一节补上。
\end{proof}

\begin{theorem}
\label{thm:y1-is-C-riemann}
$j$-不变量诱导一个双全纯同构
\[
  j\colon Y(1)=\SL_2(\Z)\backslash\HH \xrightarrow{\ \sim\ } \C.
\]
因此 $Y(1)$ 作为黎曼面就是复平面本身。
\end{theorem}

\begin{proof}[解释]
在集合层面，这件事在\cref{cor:torus-tau,sec:modular-curves-moduli} 中已经出现过：
$j(\tau)$ 恰好分类复椭圆曲线的同构类。
现在要补上的，是"它为什么与复结构相容"。

$j(\tau)$ 在 $\HH$ 上是全纯的，并且对 $\SL_2(\Z)$ 不变，
所以它下降为 $Y(1)$ 上的全纯函数。
在普通点处，这没有任何问题；
在椭圆点 $i$ 与 $\rho=e^{2\pi i/3}$ 附近，
$j$ 恰好具有与稳定子相匹配的分歧：
局部上 $j-1728$ 像二次函数，$j$ 在 $\rho$ 附近像三次函数。
也就是说，$j$ 正好把我们刚才引入的局部坐标
$u^2$、$u^3$ 吸收掉，因此在商空间上仍然是良好的坐标。

最后，$j$ 在 $Y(1)$ 上是双射，逆映射由复分析中的开映射原理自动是全纯的，
所以这是一个双全纯同构。
\end{proof}

\begin{remark}
这一节的结论可以浓缩成一句话：
\emph{模曲线不是一个粗糙的商拓扑空间，而是一个真正的复几何对象。}
从这一刻起，模形式不再只是满足变换律的函数，
而是可以被重新解释成模曲线上的截面、微分形式和除子。
这正是后面几节的主线。
\end{remark}
